\section{An odd periodic slice and its stopped geometry}
\label{sec:odd-slice}

We now work in the remaining nonisotropic case: the shape strip has finite
width \(H\), and the growth at its boundary is critical.  By
\Cref{model:angle-budget-activated}, there are constants
\(0<\rho<1\) and \(C\) such that
\begin{equation}
 \int_{k_0}^{K}\mu_f(v)^2\,\dd v\le C K^{2\rho}
 \qquad(K\ge k_0).
 \label{odd:angle-budget}
\end{equation}
We shall use this saving to control the geometry of one fixed periodic
orbit family, even when its complex continuation approaches the shape wall.
The estimates will hold up to a stopping time defined by the greatest
logarithmic shape depth visited so far. Keeping this history in the stopping rule is essential:
it permits us to use a finite collection of determinant estimates when the
sites requiring those estimates change along a path.

\subsection{The cyclic geometry and its Jacobi matrix}

Fix a sufficiently small positive parameter
\[
 t_0=\frac{2\pi}{N},\qquad N\ge3\text{ odd},
\]
inside the real analytic collar.  The integer \(N\) and the parameter \(t_0\)
are fixed throughout this section, before any limiting depth is chosen.
Set
\[
 h=h_X(t_0)=h_L(t_0),\qquad 0<h\le H,
\]
where equality of the two widths on real parameters and the upper bound
follow from the continuation results and \Cref{thm:shape-lift}.
On \(|\Im y|<h\), put
\begin{equation}
 \begin{aligned}
 s_j(y)&=\mathfrak s(y+jt_0,t_0),&
 X_j(y)&=\Gamma(s_j(y)),\\
 \psi_j(y)&=\phi(s_j(y)),&
 f_j(y)&=f(s_j(y)).
 \end{aligned}
 \label{odd:site-data}
\end{equation}
We use indices modulo \(N\), with the prescribed lifts
\(s_{j+N}=s_j+2\pi\) and \(\psi_{j+N}=\psi_j+2\pi\).
The functions \(s_j\) and \(\psi_j\) are holomorphic; the position functions
are meromorphic.

Recall
\[
 n_\theta=(\cos\theta,\sin\theta),\qquad
 e_\theta=(-\sin\theta,\cos\theta).
\]
The equations
\begin{equation}
 \theta_{j-1}+\theta_j=2\psi_j,
 \qquad \theta_{j+N}=\theta_j+2\pi
 \label{odd:angle-system}
\end{equation}
have a unique solution because \(N\) is odd.  For example, repeated
substitution expresses each \(\theta_j\) as an alternating sum of \(N\)
consecutive \(\psi\)'s plus a fixed multiple of \(\pi\).  These solutions
are holomorphic and agree, on the real collar, with the directed angles of
the billiard edges.  Consequently the identities below continue
meromorphically from that collar.  Define
\begin{equation}
 \begin{split}
 d_j^0&=\theta_j-\psi_j,\qquad
 p_j=X_j\cdot n_{\theta_j},\\
 X_{j+1}-X_j&=\ell_j e_{\theta_j},\qquad
 H_j^0=X_j\cdot n_{\psi_j}.
 \end{split}
 \label{odd:edge-data}
\end{equation}
Here and below the scalar product is the complex bilinear extension of the
Euclidean scalar product.  The chord lengths \(\ell_j\) are directed
meromorphic lengths; they are positive on the real collar.

The elementary support identities are
\begin{align}
 p_{j-1}+p_j&=2H_j^0\cos d_j^0,
 &p_j-p_{j-1}&=2(X_j\cdot e_{\psi_j})\sin d_j^0,
 \label{odd:offset-system}\\
 \dd H_j^0&=(X_j\cdot e_{\psi_j})\,\dd\psi_j,
 &\dd X_j&=c e_{\psi_j}f_j^{-2}\,\dd s_j.
 \label{odd:support-differentials}
\end{align}
Indeed, \(\theta_{j-1}=\psi_j-d_j^0\) and
\(\theta_j=\psi_j+d_j^0\), which gives the first line by addition and
subtraction.  The second line follows from
\(\Gamma'=c e_\phi f^{-2}\) and \(\phi'=f\).
Although \(X_j\) can have a pole, \(H_j^0\) is regular there.  To see the
cancellation, let \(s_0\) be a zero of \(f\).  By \Cref{lem:f-zeros},
\[
 f(s)=a(s-s_0)+O((s-s_0)^3),\qquad a\ne0.
\]
Thus \(\phi(s)-\phi(s_0)=O((s-s_0)^2)\), and the residue of \(\Gamma\)
is parallel to \(e_{\phi(s_0)}\).  Its scalar product with
\(n_{\phi(s)}\) therefore has no pole.

Oddness also makes the offset system invertible.  In fact
\[
 p_j=\sum_{m=0}^{N-1}(-1)^m H_{j-m}^0\cos d_{j-m}^0.
\]
After finding the offsets, the second equation in
\Cref{odd:offset-system} recovers the tangential component of every
\(X_j\).  In particular, if the angles and the absolute logarithms of
\(|\sin d_j^0|\) are at most \(B\), then
\begin{equation}
 \max_j|X_j|\le \exp(C(1+B))\max_j|H_j^0|.
 \label{odd:reconstruction}
\end{equation}
All constants in this section may depend on the fixed slice.

Define the symmetric matrices \(M_0\), \(\mathcal V\), and \(\mathcal H\)
by
\begin{equation}
 \langle M_0z,z\rangle
   =\sum_{j=0}^{N-1}\frac{(z_j+z_{j+1})^2}{\ell_j},\qquad
 \mathcal V_j=\frac{2f_j^3}{c\sin d_j^0},\qquad
 \mathcal H=M_0-\mathcal V,
 \label{odd:Jacobi-matrix}
\end{equation}
where \(\mathcal V\) is diagonal.  Introduce the vector
\begin{equation}
 w_j=c\sin d_j^0\,f_j^{-2}s_j'(y).
 \label{odd:w-definition}
\end{equation}

\begin{lemma}[Odd-cycle matrices]
\label{lem:odd-matrices}
The following identities hold meromorphically on the slice:
\begin{equation}
 \mathcal H w=0,\qquad
 2f_js_j'=(M_0w)_j.
 \label{odd:jacobi-angle}
\end{equation}
On the real collar, \(w_j>0\), \(\mathcal H\) is negative semidefinite,
and each of its proper principal submatrices is negative definite.
The matrix \(M_0\) is positive definite there.

More generally, let \(I\) be a proper subset of the sites and let
\(F\subseteq I^c\).  Retain the sites of \(I\) with their actual curvature,
retain the sites of \(F\) with curvature zero, and delete all other sites.
The resulting matrix
\begin{equation}
 B_{I,F}
  =(M_0)_{I\cup F,I\cup F}
    -\operatorname{diag}\bigl(\mathcal V_i\ind_{\{i\in I\}}\bigr)
 \label{odd:modified-matrix}
\end{equation}
is invertible on the real collar.  In particular its meromorphic
determinant is not identically zero.  For an empty retained set we use
determinant \(1\) and inverse norm \(0\).
\end{lemma}

\begin{proof}
First work on the real collar.  If \(\delta\mathrm{arc}_j\) is an arc-length
variation at vertex \(j\), put
\(z_j=\sin d_j^0\,\delta\mathrm{arc}_j\).  The normal projections on edge
\(j\) of the two endpoint variations are \(z_j\) and \(-z_{j+1}\).
The second variation of that edge length is therefore
\((z_j+z_{j+1})^2/\ell_j\), apart from the acceleration of the endpoint
paths.  The latter acceleration is
\(-n_{\psi_j}f_j^3/c\) per squared arc-length increment.  Pairing it with
the length gradient \(e_{\theta_{j-1}}-e_{\theta_j}\) gives
\(-2f_j^3\sin d_j^0/c\) per squared arc-length increment, or
\(-\mathcal V_j z_j^2\).  Thus \(\mathcal H\) is the length Hessian in
these variables.  Differentiation along the periodic family gives
\(\mathcal H w=0\).  Since \(\mathcal V_jw_j=2f_js_j'\), this also gives
the second identity in \Cref{odd:jacobi-angle}.  Analytic continuation
proves the meromorphic identities.

The real collar is sufficiently small that \(s_j'>0\) and
\(\sin d_j^0>0\), so \(w_j>0\).  Let \(D_w=\diag(w_j)\) and put
\(\gamma_j=w_jw_{j+1}/\ell_j>0\).  The off-diagonal entries of
\(D_w\mathcal H D_w\) are \(\gamma_j\); the identity
\(\mathcal Hw=0\) fixes its diagonal entries as
\(-\gamma_{j-1}-\gamma_j\).  Consequently
\[
 \langle D_w\mathcal H D_w z,z\rangle
   =-\sum_j\gamma_j(z_j-z_{j+1})^2.
\]
The cycle is connected.  Its only null vectors are constant, and fixing
any one coordinate to zero removes the nullspace.  This proves the two
claims concerning \(\mathcal H\).  A null vector for \(M_0\) would satisfy
\(z_{j+1}=-z_j\) for every \(j\); oddness then forces \(z=0\).

For the last assertion, arrange the retained coordinates as \(I,F\).
The truthful diagonal block is \(\mathcal H_{I,I}<0\), and the flat
diagonal block is \((M_0)_{F,F}>0\).  When both are nonempty, the Schur
complement of the flat block is
\[
 \mathcal H_{I,I}
 -(M_0)_{I,F}(M_0)_{F,F}^{-1}(M_0)_{F,I}<0.
\]
Hence the full matrix is invertible.  If one of the two blocks is empty,
one of the already proved definite-matrix statements applies.
\end{proof}

\begin{lemma}[Cofactor normalization]
\label{odd:cofactor-normalization}
Let \(A(t_0)\ne0\) be the signed area coefficient from
\Cref{eq:area-coefficient}, in the directed-edge convention used here.
Then
\begin{equation}
 w_j^2=(-1)^{N-1}\frac{A(t_0)}{N}
       \left(\prod_{i=0}^{N-1}\ell_i\right)
       \det\mathcal H_{\widehat j}.
 \label{eq:cofactor}
\end{equation}
Here \(\mathcal H_{\widehat j}\) denotes the principal matrix with site
\(j\) deleted.
\end{lemma}

\begin{proof}
On the collar let \(v_j\) be the variation obtained by differentiating
\(\mathfrak s(y+jt,t)\) with respect to \(t\), including its argument
\(jt\), and multiplying by \(c\sin d_j^0 f_j^{-2}\), at \(t=t_0\).
Periodicity gives
\[
 v_{j+N}=v_j+Nw_j.
\]
The endpoint normal differentials on edge \(j\) are
\[
 n_{\theta_j}\cdot\dd X_j=w_j\,\dd y+v_j\,\dd t,
 \qquad
 n_{\theta_j}\cdot\dd X_{j+1}
       =-w_{j+1}\,\dd y-v_{j+1}\,\dd t.
\]
Differentiating \(X_{j+1}-X_j=\ell_je_{\theta_j}\) and taking the normal
component yields
\[
 \dd\theta_j
   =\frac{(w_j+w_{j+1})\,\dd y+(v_j+v_{j+1})\,\dd t}{\ell_j}.
\]
Since
\(\dd p_j\wedge\dd\theta_j
 =(n_{\theta_j}\cdot\dd X_j)\wedge\dd\theta_j\), the area identity gives
\[
 \frac{w_jv_{j+1}-v_jw_{j+1}}{\ell_j}=A(t_0).
\]
The shift of the line phase by \((j+1/2)t\) does not alter the coefficient
of \(\dd y\wedge\dd t\).  Summing the successive differences of
\(v_j/w_j\) therefore gives
\begin{equation}
 \sum_j\frac{\ell_j}{w_jw_{j+1}}=\frac{N}{A(t_0)}.
 \label{odd:Wronskian-sum}
\end{equation}

The deleted cofactor of the negative weighted cycle Laplacian is
\[
 (-1)^{N-1}\sum_k\prod_{i\ne k}\gamma_i
   =(-1)^{N-1}\left(\prod_i\gamma_i\right)
                         \left(\sum_i\gamma_i^{-1}\right).
\]
For completeness, write the positive Laplacian as an incidence matrix,
times \(\diag(\gamma_i)\), times its transpose, and delete one vertex.
Cauchy--Binet expands the determinant over the omitted edge.  Each
remaining edge set is a path tree, whose incidence minor has determinant
\(\pm1\); this gives precisely the displayed sum.  Undoing the conjugation
by \(D_w\), and then using \Cref{odd:Wronskian-sum}, proves
\Cref{eq:cofactor} on the collar.  Both sides are meromorphic functions of
\(y\), so the identity holds throughout the slice.
\end{proof}

\subsection{Contact with the shape wall}

\begin{theorem}[Boundary contact]
\label{thm:boundary-contact}
At almost every point of either horizontal boundary of the phase strip
\(|\Im y|<h\), the functions \(s_0,\ldots,s_{N-1}\) have finite
nontangential limits, and at least one of these limits lies on
\(|\Im s|=H\).
\end{theorem}

\begin{proof}
By \Cref{thm:shape-lift}, each \(e^{\ii s_j}\) is bounded and bounded away
from zero.  The boundary limit theorem for bounded holomorphic functions
gives nontangential limits almost everywhere
\cite[Chapter~I, Theorem~5.3]{Garnett2007}.  Near a nonzero limiting
value, choose a local logarithm.  On a connected truncated approach cone
it differs from \(\ii s_j\) by one fixed integral multiple of \(2\pi\ii\).
Thus \(s_j\) itself has a finite nontangential limit.

Suppose the limiting states are all strictly inside the shape strip on a
set of positive measure.  At a point of this set we can require, in
addition, that every \(f_j\), \(\sin d_j^0\), \(\ell_j\), and
\(\det\mathcal H_{\widehat j}\) have a nonzero limit.  To justify the
requirement, first exclude limiting zeros of the \(f_j\).  Each \(f_j\)
is holomorphic on the phase strip and nonzero on the real collar.  If its
angular limit were zero on a positive-measure subset, the same boundary
uniqueness theorem used below would force \(f_j\equiv0\), a contradiction.
Thus the finite union of these exceptional sets has measure zero.
The sine and length evaluations are then regular at the interior limiting
state tuple. Exclude their zero limits next; all cofactor evaluations are
regular after those exclusions. If any one of these meromorphic functions
had limiting value zero
on a positive-measure set, Privalov's nontangential boundary uniqueness
theorem for meromorphic functions would make it identically zero
\cite[Chapter~II, p.~91]{Garnett2007}.  This is impossible by its
nondegenerate real values.  There are only finitely many functions to
exclude.  To apply the disk boundary theorems to this strip, restrict to a
compact boundary subarc of positive measure and use a conformal map to the
disk.  The map extends analytically across that subarc with nonzero
derivative, preserving null sets and nontangential approach regions.

Write \(\mathbf s=(s_0,\ldots,s_{N-1})\).  Near such a limiting tuple the
cofactor identity and \Cref{odd:w-definition} express
\[
 (s_j')^2=G_j(\mathbf s),
\]
where each \(G_j\) is analytic and nonzero.  Choose the square roots with
the signs of the actual functions on a truncated approach cone.  We obtain
a regular analytic vector equation
\(\mathbf s'=\mathbf g(\mathbf s)\).  Local solutions starting sufficiently
near the limiting tuple exist in disks of a common positive radius.
Starting at an interior point sufficiently close to the boundary, uniqueness
identifies that solution with the actual slice on their overlap.  This
continues the whole slice arc across the chosen boundary point.

It remains to explain why an extension of this single slice yields a
\emph{joint} extension, as required for \Cref{thm:no-patch}.  We give the
argument at the upper boundary.  The states
\[
 Z=(s_0,s_1)
\]
on the extended arc form a regular analytic curve, parameterized by \(y\).
All incidence angles and chords along its orbit are nondegenerate.  Thus
the local billiard return \(R\), after \(N\) reflections and subtraction
of the known \(2\pi\) shifts, is analytic near these states.  At each
reflection, this follows by solving the next incidence equation on the
regular table arc; its derivative in the next intersection coordinate is
nonzero because the sine of incidence is nonzero.  The return fixes the
extended curve, by analytic identity from below.

The determinant of its linearization at a fixed state is one.  One can
see this without choosing symplectic coordinates: successive linearized
stationarity equations in the Jacobi variables have transfer determinants
equal to ratios of consecutive off-diagonal entries.  The ratios telescope
over the cycle.  The coordinate factors converting \(\delta s_j\) to
\(c\sin d_j^0 f_j^{-2}\delta s_j\) are also periodic along the fixed orbit.
Choose local coordinates \((u,v)\) with the fixed arc given by
\((u,v)=(y,0)\).  The return consequently has the form
\begin{equation}
 R(u,v)=
 \bigl(u+\sigma(u)v+O(v^2),\ v+\tau(u)v^2+O(v^3)\bigr).
 \label{odd:return-expansion}
\end{equation}

In a lower part of this neighborhood the original joint data are known.
For
\[
 Z(y,t)=\bigl(\mathfrak s(y,t),\mathfrak s(y+t,t)\bigr),
 \qquad b=N(t-t_0),
\]
the reflection identities give
\(R(Z(y,t))=Z(y+b,t)\).  Expanding there in the coordinates just chosen,
write
\[
 u=y+O(b),\qquad v=b p_0(y)+O(b^2).
\]
The first component of the return identity at order \(b\), and its second
component at order \(b^2\), give respectively
\begin{equation}
 \sigma p_0=1,\qquad p_0'=\tau p_0^2,
 \qquad \sigma'=-\tau.
 \label{odd:return-compatibility}
\end{equation}
These identities continue along the extended arc.  By moving the boundary
point slightly within this arc if necessary, take a smaller neighborhood
where \(\sigma\ne0\), and set \(p_0=1/\sigma\).

For small nonzero \(b\) with \(|\Im b|\ge |b|/3\), the return and its
inverse continue the joint states upwards in steps \(\pm b\).  Here are
the relevant uniform estimates.  In the coordinates \((u,p)=(u,v/b)\),
\Cref{odd:return-expansion} reads
\[
 (u,p)\longmapsto(u,p)
       +b\bigl(\sigma(u)p,\tau(u)p^2\bigr)+O(b^2)
\]
on fixed compact subsets; the inverse has the negative first increment.
The analytic differential equation with that vector field has the regular
solution \((u,p)=(y,p_0(y))\), by
\Cref{odd:return-compatibility}.  Choose the arc neighborhood and the
height to be crossed sufficiently small.  The steepness condition keeps
all the slanted paths from the lower starting slab inside that neighborhood.
There are \(O(1/|b|)\) steps; each has error \(O(|b|^2)\) and stability
factor at most \(1+C|b|\).  The continued states therefore stay within
\(O(|b|)\) of this regular solution.  All shape coordinates remain strictly
inside the shape strip.

All intermediate site coordinates are recovered by the regular analytic
billiard steps along the orbit.  Thus their continuations agree with the
original \(s_j\)'s wherever the latter were already defined.

Use locally constant numbers of backward steps to place initial states
in the lower slab.  On overlaps, the exact return identity in that slab
and uniqueness of each forward step identify the constructions.  They
give joint holomorphic germs in the steep sectors, attached by vertical
continuation to the lower data.

Finally choose a small \(\eta>0\), then a much smaller \(r>0\), and use
the analytic disks
\begin{equation}
 t=t_0+r\zeta,\qquad
 y=x_*+\ii\varsigma-\ii\eta\zeta^2+w,
 \qquad |\zeta|\le1,
 \label{odd:contact-disks}
\end{equation}
as \(\varsigma\) rises from \(h-2\eta\) to \(h\); \(w\) varies in a
sufficiently small complex disk.  On the boundary arcs where \(\zeta\)
is near the real axis, the phase point is a fixed amount below the wall.
For small \(r\), these arcs lie in the original joint tube, by lower
semicontinuity of the width.  On the remaining arcs the parameter is in
a fixed steep sector, so the preceding continuation applies.  The collar
germs agree because both are vertical continuations of the lower data.
The full disks are initially inside known data.  Holomorphic product
continuation, as in \Cref{lem:meromorphic-hartogs}, fills them as the center
is raised.  Equivalently, pull the germs back to a collar of the unit circle
and a complex neighborhood of the center path: the negative Cauchy moments
vanish in the initial open part and hence throughout that connected
neighborhood. More explicitly, combine \(x_*+\ii\varsigma+w\) into
one complex center variable \(a\). The map
\((a,\zeta)\mapsto(a-\ii\eta\zeta^2,t_0+r\zeta)\)
has Jacobian determinant \(r\ne0\), so the filling supplies a joint
neighborhood in \((y,t)\).

The maximum principle preserves \(|\Im s_j|<H\) on the filled disks.
Composition with \(\Gamma\) gives a joint meromorphic patch for \(X\) at
\((x_*+\ii h,t_0)\), agreeing with the original tube below the wall.
This contradicts \Cref{thm:no-patch}, which applies also to the rational
parameter \(t_0\).  The lower boundary is treated in the same way, or by
real symmetry.  The supposed positive-measure set therefore cannot exist.
\end{proof}

\subsection{Brownian motion, history levels, and fixed stopping rules}

We use the standard conformal time-change, optional-stopping, and
martingale maximal theorems for Brownian motion; see
\cite[Chapters~II, IV and V]{RevuzYor1999}. Their applications, including the stopping cuts,
are specified below.
Let \(y_u\) be planar Brownian motion, started with \(y_0\) uniform on one
real period, and let \(\zeta\) be its first exit from \(|\Im y|<h\).
All functions are evaluated on the real lift of this path.  Put
\begin{equation}
 \begin{aligned}
 \delta_j(u)&=H-|\Im s_j(y_u)|,&
 k_j(u)&=-\log\delta_j(u),&
 K_j(u)&=\max_{0\le v\le u}k_j(v),\\
 r_y(u)&=h-|\Im y_u|,&
 n_y(u)&=-\log r_y(u),&
 \tau_L&=\inf\{u:\max_jK_j(u)=L\}.
 \end{aligned}
 \label{odd:stopping-data}
\end{equation}
For sufficiently large \(L\), \(\tau_L<\zeta\) almost surely.  Indeed the
bounded functions \(e^{\ii s_j}\) are Poisson integrals of their almost
everywhere boundary values.  This follows first on narrower strips and
then on the whole strip by dominated convergence.  Along Brownian motion
they are the conditional-expectation martingales of their terminal boundary
values.  Announcing the exit by exits from narrower strips shows that their
limits at Brownian exit are those terminal values.  The exit distribution
is absolutely continuous on the horizontal boundary, so
\Cref{thm:boundary-contact} makes at least one \(\delta_j\) tend to zero.
Continuity then gives every sufficiently large history level before exit.

The isolated zeros and poles of the meromorphic slice functions are polar
for planar Brownian motion and are almost surely avoided.  All statements
involving them can first be stopped on compact subsets avoiding these
points, followed by an exhaustion.  In what follows, ``with probability
tending to one through \(\tau_L\)'' means that one event of probability
\(1-o(1)\) supports the stated bounds simultaneously at every time
\(0\le u\le\tau_L\) and at every site.

\begin{lemma}[Shape time change and its consequences]
\label{odd:time-change}
Each process \(s_j(y_u)\), before phase exit, is a Brownian motion in the
shape strip under the scalar time change
\[
 Q_j(u)=\int_0^u|s_j'(y_v)|^2\,\dd v,
\]
possibly stopped before shape exit.  Its initial horizontal law has a
density bounded by a fixed constant relative to uniform measure modulo
\(2\pi\).  In particular:
\begin{enumerate}[label=\textup{(\roman*)}]
 \item The expected occupation of any nonnegative periodic shape
 observable, integrated against \(\dd Q_j\), is bounded by a fixed
 constant times its ordinary shape-strip occupation from a uniform real
 start.  The total shape clocks have bounded expectation.  Moreover,
 \begin{equation}
  \E\int_0^{\tau_L}
       \ind_{\{k_j(u)\le K\}}|f_j(y_u)s_j'(y_u)|^2\,\dd u
       \le C(1+K)^{2\rho}
  \label{odd:angle-clock}
 \end{equation}
 for large \(K\).
 \item At the first hit of a deterministic shape-depth row, the
 subprobability law of the hit is dominated by a fixed constant times the
 corresponding law for ordinary shape Brownian motion from a uniform real
 start.  In particular that ordinary horizontal law is uniform.
 \item After a first hit of shape depth \(\delta\), the remaining shape
 clock before any later phase stop has tail at most
 \begin{equation}
  C\min\left(1,\frac{\delta}{\sqrt T}\right)
  \label{odd:remaining-shape-clock}
 \end{equation}
 at clock time \(T>0\).  This estimate also holds conditionally at that
 hit whenever it occurs.
 \item For a sufficiently large fixed \(C\), with probability tending
 to one through \(\tau_L\),
 \begin{equation}
  k_j(u)\ge K_j(u)-C\log L\qquad\text{for every }j.
  \label{odd:retreat}
 \end{equation}
\end{enumerate}
\end{lemma}

\begin{proof}
The real and imaginary parts of a holomorphic function of planar Brownian
motion are local martingales.  The Cauchy--Riemann equations give equal
quadratic variations \(\dd Q_j\) and zero mutual quadratic covariation.
The Brownian time-change theorem therefore applies.  To use it up to a
given interior phase stop, append independent Brownian increments in
intrinsic time after that stop.  The extension cannot have exited the
shape strip before the part of the path already used, because all those
states lie in the strip.  Approximation by compact interior stops proves
the same domination for nonnegative occupations over the whole phase
interval.

On the real collar \(y\mapsto s_j(y)\) is an increasing analytic circle
diffeomorphism whose derivative is bounded above and below.  This proves
the initial-density assertion.  Under a uniform real start, the periodic
shape-strip Green weight is a constant times
\(\delta=H-|\Im s|\).  It has finite integral, proving the total-clock
claim.  At either wall, the contribution of \(|f|^2\) up to log depth
\(K\) is bounded by a constant times
\[
 \int_{e^{-K}}^{\delta_*}\delta\,\E|f|^2\,\dd\delta
   =\int_{k_*}^{K}\delta^2\E|f|^2\,\dd k
   =\int_{k_*}^{K}\mu_f(k)^2\,\dd k.
\]
The fixed interior region has bounded contribution.  Now use
\Cref{odd:angle-budget}.  Domination of first-hit laws follows from the
same Brownian extension; if the original path never reaches the row,
it contributes no mass.

Starting at distance \(\delta\) from one horizontal wall, the exit time
from the strip is no greater than the first time the vertical Brownian
coordinate hits that particular wall.  The one-dimensional reflection
principle bounds survival to time \(T\) by
\(C\min(1,\delta/\sqrt T)\).  Applying this after the hit, with the
strong Markov property in intrinsic time, proves (iii).

For (iv), at the first hit of each integer level \(m\), the probability
of reaching depth \(R e^{-m}\) before hitting the adjacent wall is at
most \(1/R\), by the linear one-dimensional hitting probability.  Take
\(R\) a sufficiently large fixed power of \(L\), and sum over the
\(O(L)\) integer levels up to \(L\), and the fixed number of sites.
If a retreat larger than \(C\log L\) occurred, it would contain one of
these integer-level retreats, with an immaterial factor from rounding.
The union probability tends to zero.
\end{proof}

We record explicitly the two localization principles used below.

\begin{lemma}[Logarithmic deficit at a stopping cut]
\label{odd:log-deficit}
Let \(g\) be a nontrivial meromorphic slice function, with
\(\log|g(y_0)|\ge-C_0\) on the real starting row.  Suppose
\(\log|g(y_u)|\le B_0\), where \(B_0\ge1\), up to a specified stopping
cut \(T\).  Then
\begin{equation}
 \Prob\left\{\inf_{u\le T}\log|g(y_u)|<-B_1\right\}
     \le \frac{B_0+C_0}{B_0+B_1}
     \le\frac{C(1+B_0)}{1+B_1}
 \label{odd:log-deficit-probability}
\end{equation}
for \(B_1\ge1\).  Earlier bounds may be incorporated by stopping at their
first failure; their failure probabilities are then added separately.
\end{lemma}

\begin{proof}
Away from the polar zero and pole set, \(\log|g(y_u)|\) is a continuous
local martingale.  Before the indicated cut,
\(B_0-\log|g(y_u)|\) is nonnegative, hence is a supermartingale after
localization.  Stop additionally at its first hit of \(B_0+B_1\).
The supermartingale inequality gives the first bound.  Exhaustion removes
the auxiliary compact localizations.
\end{proof}

For an analytic differential integrated after a stopping time, we shall
likewise apply the maximal inequality only after a quadratic-time cut.
More precisely, if its scalar quadratic time is at most a deterministic
\(B\) on some collection of paths, stop the integral at quadratic time
\(2B\) before estimating its maximum.  The resulting second-moment bound
is \(O(B)\).  It applies to the original integral on those paths, without
conditioning on the collection of paths.  This remains true when the
integrand is masked by adapted status indicators, using their predictable
versions.  All ensuing probability estimates use these cuts.

Put
\begin{equation}
 a_j=\log|f_j|,
 \qquad
 \mathcal G(u)=\max_j\left(
  |\psi_j|+\log(1+|X_j|)
       +|\log|\sin d_j^0||+|\log|\ell_j||\right)(y_u).
 \label{odd:geometry-size}
\end{equation}
By \Cref{cor:critical-polynomial-bounds}, everywhere before the relevant
history stop,
\begin{equation}
 a_j\le k_j^++C\log(2+k_j^+)+C,
 \qquad k_j^+=\max(k_j,0).
 \label{odd:amplitude-upper}
\end{equation}

\begin{lemma}[A finite family of modified inverse estimates]
\label{odd:modified-inverses}
Let \(V\ge U\ge1\) be deterministic.  For each fixed proper truthful set
\(I\) and flat set \(F\subseteq I^c\), stop at \(\tau_L\), at the first
failure of \(\mathcal G\le U\), or at the first time an \(I\)-history
exceeds \(U\), whichever comes first.  Except on an event of probability
at most
\begin{equation}
 \frac{C(1+U)}{1+V},
 \label{odd:inverse-failure}
\end{equation}
the bound
\begin{equation}
 \|B_{I,F}^{-1}\|\le\exp(CV)
 \label{odd:inverse-bound}
\end{equation}
holds throughout this stopped interval.  The constant is uniform over
the finitely many choices of \(I,F\).  Thus one may secure all these
bounds by a finite union, and subsequently choose \(I,F\) as functions of
the current state before the common geometry cut, whenever their selected
history caps hold.
\end{lemma}

\begin{proof}
On the stopped interval, \Cref{odd:amplitude-upper} bounds \(a_i\) by
\(C(1+U)\) for \(i\in I\).  The geometry cap bounds the reciprocal
lengths, the reciprocal incidence sines, and the angles.  Hence every
entry, determinant, and cofactor of \(B_{I,F}\) has modulus at most
\(\exp(C(1+U))\).  There is no amplitude bound required at flat or deleted
sites: their curvature terms do not appear.  By \Cref{lem:odd-matrices},
the determinant has a positive minimum in modulus on the compact real
starting period.  Apply \Cref{odd:log-deficit} to that determinant with
upper bound \(C(1+U)\) and lower-log threshold \(-V\).  The adjugate
formula gives \Cref{odd:inverse-bound} outside the resulting exceptional
event.  Empty matrices require no estimate.

There are at most \(3^N\) choices.  Each matrix is defined and estimated
from the original real start, with its own fixed stopping rule.  If at
a later time all the histories selected by that matrix are at most \(U\),
none of them has crossed its cap previously.  Thus the estimate already
secured for that fixed candidate is still valid as long as the common
geometry cut has not occurred.  This is the reason for
using history levels rather than current depths.  No estimate is restarted
when the current selection changes.
\end{proof}

\subsection{An initial bound below linear order}

\begin{lemma}[Initial stopped geometry]
\label{odd:initial-geometry}
There is a fixed \(\sigma<1\) such that
\[
 \mathcal G(u)\le L^\sigma\qquad(0\le u\le\tau_L)
\]
with probability tending to one as \(L\to\infty\).
\end{lemma}

\begin{proof}
The angle differential is \(\dd\psi_j=f_j\,\dd s_j\).  Up to \(\tau_L\)
all current log depths are at most \(L\), so
\Cref{odd:angle-clock} and the maximal inequality bound all angles by
\(L^{\rho+\eps}\) with probability tending to one, for each fixed
\(\eps>0\).  The linear equations \eqref{odd:angle-system} give the same
bound for \(\theta_j,d_j^0\).  They imply an upper bound of this order for
\(\log|\sin d_j^0|\).  The sine has a positive minimum on the real
starting collar.  Applying \Cref{odd:log-deficit} with a slightly larger
power gives, for some fixed \(\rho<\rho'<1\), simultaneous angle and sine
cost at most
\[
 r_*=L^{\rho'}.
\]
We stop at the first failure of these bounds for the estimates that follow,
and add the already vanishing failure probability at the end.

Choose
\[
 G=L^\gamma,\qquad \frac{1+\rho'}2<\gamma<1,
\]
and a small fixed \(b>0\).  At each site use two statuses, initially
moderate.  Switch from moderate to tiny at a downcrossing of
\(a_j=-2bG\), and switch back at an upcrossing of \(a_j=-bG\).
The hysteresis gap is \(bG\).  By \Cref{odd:amplitude-upper},
\(CL-a_j\) is a nonnegative local martingale, hence a supermartingale,
up to \(\tau_L\) and the preceding cuts. Its upcrossing levels are
\(CL+bG\) and \(CL+2bG\). The upcrossing inequality, and alternation
of the two switch directions, give an expected number of switches
\(O(L/G)\). One can also obtain this
bound by holding supermartingale increments during attempted passages
between the two levels: each completed passage earns \(bG\), while an
unfinished passage has possible loss bounded by \(CL+2bG\).
Thus, for any sufficiently small fixed \(\eps_1>0\), the total number of
switches over the fixed number of sites is at most
\begin{equation}
 L^{1+\eps_1}/G
 \label{odd:switch-count}
\end{equation}
with probability tending to one.

Let
\[
 \mathcal R(u)=1+\sup_{0\le v\le u,\,j}|X_j(y_v)|.
\]
During a moderate interval, \(a_j\ge-2bG\).  From
\Cref{odd:support-differentials}, the coefficient of \(\dd s_j\) in
\(\dd X_j\) has modulus at most \(e^{5bG}\), for large \(L\), because
\(r_*=o(G)\).  During a tiny interval, \(a_j\le-bG\), and the coefficient
of \(\dd s_j\) in \(\mathcal R^{-1}\dd H_j^0\) has modulus at most
\(e^{-bG/2}\).  The latter estimate uses
\[
 \mathcal R^{-1}|X_j\cdot e_{\psi_j}|\,|f_j|
       \le \exp(Cr_*-bG).
\]
Mask these two differentials by their corresponding status indicators.
The total expected shape clocks are bounded by
\Cref{odd:time-change}.  The maximal inequality, with the cuts described
above, shows that the cumulative moderate-position integrals and
tiny-support integrals have respective suprema at most
\begin{equation}
 e^{6bG},\qquad e^{-bG/3}
 \label{odd:masked-integrals}
\end{equation}
with probability tending to one.

On a tiny interval, write \(\dd H_j^0=\mathcal R\,\dd M_j\), where the
appropriate masked martingale \(M_j\) is bounded as in
\Cref{odd:masked-integrals}.  Integration by parts against the continuous
increasing function \(\mathcal R\) bounds the support increment on any
such interval by
\(C e^{-bG/3}\mathcal R\) at its endpoint.  For example, after subtracting
the starting value of \(M_j\), both the endpoint term and the integral
against \(\dd\mathcal R\) are bounded by that quantity.

Between successive switches of the joint status vector, a moderate site's
position differs from its starting position by at most \(2e^{6bG}\), and
a tiny site's support has the preceding increment bound.  If
\(\mathcal R_{\rm start}\) is the running maximum at the start of such
an interval, every support during it is therefore at most
\begin{equation}
 e^{Cr_*}\bigl(\mathcal R_{\rm start}+C e^{6bG}\bigr)
       +C e^{-bG/3}\mathcal R.
 \label{odd:support-interval-bound}
\end{equation}
Use \Cref{odd:reconstruction} and absorb the last term; its coefficient
after reconstruction tends to zero because \(r_*=o(G)\).  Each interval
then multiplies the maximum of its starting bound and \(C e^{6bG}\) by
at most \(e^{Cr_*}\).  By \Cref{odd:switch-count}, its accumulated
logarithmic cost is at most
\[
 C r_*L^{1+\eps_1}/G=o(G),
\]
provided \(\eps_1\) is small enough that
\(2\gamma>1+\rho'+\eps_1\).  Hence \(\log\mathcal R\le C G\) with
probability tending to one.

The identity \(\ell_j=(X_{j+1}-X_j)\cdot e_{\theta_j}\) now gives upper
logarithmic length bounds of the same order.  Apply
\Cref{odd:log-deficit} to each nonzero length, allowing a further small
increase of the exponent, to obtain lower logarithmic bounds.  Choose a
fixed \(\sigma\) strictly between all these exponents and \(1\).
Constants are then absorbed by \(L^\sigma\) for large \(L\), proving the
claim.  Each previous good bound was enforced by its first-failure cut;
only the sum of their vanishing exceptional probabilities has been added.
\end{proof}

\begin{lemma}[Phase depth and phase clock]
\label{odd:phase-clock}
For some fixed \(C\), with probability tending to one through \(\tau_L\),
\begin{equation}
 n_y(u)\le CL,\qquad
 \int_0^{\tau_L}r_y(u)^{-2}\,\dd u\le L^C.
 \label{odd:phase-clock-bounds}
\end{equation}
\end{lemma}

\begin{proof}
Apply \Cref{odd:initial-geometry} through \(\tau_{2L}\).  The entries of
\(\mathcal H\) have upper logarithmic size \(O(L)\), by
\Cref{odd:amplitude-upper}; the incidence sines and lengths have only
sublinear logarithmic cost.  The cofactor formula
\Cref{eq:cofactor} consequently gives
\begin{equation}
 |w_j|+|s_j'|\le e^{C_1L}
 \label{odd:rough-derivative-bound}
\end{equation}
through that stop on the good paths.  The second bound follows from
\Cref{odd:w-definition}, using the amplitude upper bound again.

Let \(T\) be the first phase depth \(r_y=e^{-C_2L}\), with \(C_2\) fixed
and sufficiently large.  By the one-dimensional hitting estimate, after
\(T\) the remaining phase time is at most \(e^{-C_2L}\), except with
probability at most \(C e^{-C_2L/2}\).  Suppose \(T<\tau_L\).  At time
\(T\), all shape depths are at least \(e^{-L}\), whereas at
\(\tau_{2L}\) one of them is \(e^{-2L}\).  On the good derivative paths,
the quadratic time of each shape increment from \(T\) to
\(\tau_{2L}\) is bounded by
\(e^{2C_1L-C_2L}\).  Enforce this by the derivative cut and by a physical
time cut of \(e^{-C_2L}\) after \(T\).  The maximal inequality then shows
that, for sufficiently large \(C_2\), all these increments are
\(o(e^{-L})\) with probability tending to one.  This contradicts the
required change of one imaginary part by at least
\(e^{-L}-e^{-2L}\).  Thus \(T<\tau_L\) has probability tending to zero,
which proves the first bound.

Under the uniform real start, the phase-strip Green weight at a wall is
proportional to \(r_y\).  Therefore
\[
 \E\int_0^\zeta
    \ind_{\{r_y\ge e^{-C_2L}\}}r_y^{-2}\,\dd u
       \le C(1+L).
\]
On the already established good depth event, the integral through
\(\tau_L\) is part of this truncated occupation.  Markov's inequality,
for example at threshold \(L^2\), gives a polynomial bound with
probability tending to one.  Increasing the single constant \(C\) if
necessary yields \Cref{odd:phase-clock-bounds}.
\end{proof}

\subsection{Improvement to every positive power}

\begin{theorem}[Stopped geometry at every power]
\label{thm:stopped-geometry}
For every fixed \(\sigma>0\), with probability tending to one as
\(L\to\infty\),
\begin{equation}
 \mathcal G(u)\le L^\sigma\qquad(0\le u\le\tau_L).
 \label{odd:all-power-geometry}
\end{equation}
The phase-depth and phase-clock bounds of \Cref{odd:phase-clock} hold
simultaneously.  All estimates concern the same fixed odd slice.
\end{theorem}

\begin{proof}
It suffices to improve an exponent \(0<\sigma<1\) already known to work.
Choose a small fixed \(\eps>0\) with \(\sigma+\eps<1\), and put
\[
 H_0=L^{\sigma+\eps}.
\]
The old estimate, applied at level \(H_0\), gives geometry at most
\(H_0^\sigma\) through \(\tau_{H_0}\), with probability tending to one.
We shall show that, after this time, angles and supports change by amounts
exponentially smaller than this improved initial scale.

Choose a sufficiently small fixed \(\eta>0\), and a finite family of
disjoint intervals \((\xi_m,\xi_m+\eta)\) inside
\((\sigma+\eps/3,\sigma+\eps)\), with more than \(N\) members.  At every
time after \(\tau_{H_0}\), at least one history is \(\ge H_0\).  There is
one choice \(\xi=\xi_m\) for which the corresponding power interval
\((L^\xi,L^{\xi+\eta})\) contains none of the \(N\) history levels.
For that choice, partition the sites as
\[
 I=\{j:K_j\le L^\xi\},\qquad
 I^c=\{j:K_j\ge H_*\},\qquad H_*=L^{\xi+\eta}.
\]
Endpoint equalities can be assigned to their respective sets; they do not
affect any estimate.  The set \(I\) is proper because one history is
\(\ge H_0\).  By \Cref{odd:retreat}, every deep site's current log depth
satisfies \(k_j\ge H_*/2\), on an event of probability tending to one.

We secure all matrix inverses before making any state-dependent selection.
For each fixed \(\xi_m\), proper set \(I\), and flat subset \(F\subseteq
I^c\), use \Cref{odd:modified-inverses} with, for instance,
\[
 U=L^{\xi_m+\eta/4},\qquad V=L^{\xi_m+\eta/2}.
\]
The old geometry cap \(L^\sigma\) and the selected truthful history cap
\(L^{\xi_m}\) both lie below \(U\) for large \(L\).  The failure probability
is \(O(L^{-\eta/4})\) for each candidate.  The number of candidates is
fixed, so outside one event of probability tending to zero, all required
inverses have norm
\begin{equation}
 \|B_{I,F}^{-1}\|\le
       \exp\bigl(C L^{\xi+\eta/2}\bigr)
       =\exp(o(H_*)).
 \label{odd:gap-inverse}
\end{equation}
These bounds hold whenever the corresponding histories are below their
caps, regardless of how often the empty gap and the flat subset change.

Set
\[
 z_j=r_y w_j.
\]
Schwarz--Pick, transported from the disk by conformal maps of the two
strips, gives \(r_y|s_j'|\le C\delta_j\)
\cite[Chapter~I, Lemma~1.2]{Garnett2007}; the strip hyperbolic densities
are comparable with reciprocal boundary distance.  Hence
\Cref{odd:w-definition} and the old geometry bound imply
\begin{equation}
 \log|z_j|+2a_j+k_j\le C L^\sigma,
 \label{odd:Schwarz-suppression}
\end{equation}
with the usual interpretation when \(z_j=0\).
Fix a sufficiently small positive constant \(c_1\).  Among the deep sites
retain as flat exactly those for which
\[
 |z_j|>e^{-c_1H_*};
\]
call this set \(F\), and delete the other deep sites.  At a retained deep
site, \Cref{odd:Schwarz-suppression} and \(k_j\ge H_*/2\) imply
\(a_j<-c_2H_*\) for a fixed \(c_2>0\), since \(L^\sigma=o(H_*)\).
Consequently its actual curvature term obeys
\[
 |\mathcal V_j|\le e^{-c_3H_*}
\]
for another fixed positive constant.  Every deleted coordinate already
has modulus at most \(e^{-c_1H_*}\).

Apply the equation \((M_0-\mathcal V)z=0\) to the retained coordinates
\(I\cup F\).  Its left matrix is \(B_{I,F}\) minus the tiny curvature
diagonal on \(F\); its right side is the coupling to the deleted
coordinates.  Entries of \(M_0\) cost at most \(e^{CL^\sigma}\).
By \Cref{odd:gap-inverse}, the tiny diagonal is absorbable, and the
right side gives
\begin{equation}
 \max_j|z_j|\le e^{-c_4H_*}
 \label{odd:all-sites-suppressed}
\end{equation}
for a fixed \(c_4>0\), uniformly over all the finite choices.  The same
argument includes an empty deleted set, in which case the retained
homogeneous system has only the zero solution.

The important consequence is suppression at \emph{every} site, including
the shallow ones.  From \Cref{odd:jacobi-angle},
\begin{equation}
 2r_y\partial_y\psi_j
      =\mathcal V_j z_j=(M_0z)_j.
 \label{odd:all-angle-suppression}
\end{equation}
Since \(H_*\ge L^{\sigma+\eps/3}\), the right side has modulus at most
\(\exp(-cL^{\sigma+\eps/3})\), with a fixed \(c>0\), after decreasing
that constant if necessary.  The phase-clock bound now makes the quadratic
time of each angle increment after \(\tau_{H_0}\) exponentially small.
Cut at that deterministic quadratic-time bound and apply the maximal
inequality.  We obtain, with probability tending to one,
\begin{equation}
 \sup_{\tau_{H_0}\le u\le\tau_L}
   |\psi_j(y_u)-\psi_j(y_{\tau_{H_0}})|
        \le \exp(-c' L^{\sigma+\eps/3}).
 \label{odd:late-angle-increments}
\end{equation}
All choices of history gaps and flat sets have disappeared from this
martingale statement: they were used only to prove the uniform pointwise
differential bound on the fixed good event.

The angles at \(\tau_{H_0}\) have size at most \(H_0^\sigma\).
Their incidence sines have modulus at least \(e^{-H_0^\sigma}\).
The absolute change of each sine is at most the angle change times
\(\exp(CL^\sigma)\), using the old angle bound along the interval.
By \Cref{odd:late-angle-increments}, this is negligible relative to its
improved initial lower bound.  Thus the angles and both logarithmic sine
bounds keep cost \(O(H_0^\sigma)\) throughout the late interval.

The supports obey the same improvement.  By
\Cref{odd:support-differentials},
\[
 \mathcal R^{-1}\dd H_j^0
    =\mathcal R^{-1}(X_j\cdot e_{\psi_j})\,\dd\psi_j.
\]
Its coefficient relative to the angle differential has modulus at most
\(e^{CL^\sigma}\) by the old bound.  Combining
\Cref{odd:all-angle-suppression} with the phase clock shows that the
martingale integral of this differential over the single late interval
has exponentially small supremum.  Integration by parts against the
increasing \(\mathcal R\), as in the initial geometry argument, gives
\[
 |H_j^0(y_u)-H_j^0(y_{\tau_{H_0}})|
    \le e^{-c''L^{\sigma+\eps/3}}\mathcal R(u)
\]
on a further event of probability tending to one.  Reconstruct positions
with \Cref{odd:reconstruction}, now using the improved angle and sine cost
\(O(H_0^\sigma)\).  The initial positions and supports have this same
logarithmic cost.  The coefficient of \(\mathcal R(u)\) in the last
display remains exponentially small after reconstruction and is absorbed.
It follows that
\[
 \log\mathcal R(u)\le C H_0^\sigma
       =C L^{\sigma(\sigma+\eps)}
       \qquad(u\le\tau_L)
\]
with probability tending to one.

The length upper logs have the same improved bound.  For their lower logs,
cut at the first failure of this upper bound. Choose \(\eps\) small
enough that
\[
\sigma(\sigma+\eps)<\frac{\sigma+\sigma^2}{2}<\sigma.
\]
Choose a power strictly between the first two quantities as the lower
threshold in \Cref{odd:log-deficit}; its failure probability tends to zero.
All fixed constants and all needed lower-log slack are then absorbed by
the new exponent \((\sigma+\sigma^2)/2\).  This proves an improvement of
the old assertion.

Starting with \Cref{odd:initial-geometry}, iterate this improvement a
finite number of times.  The sequence generated by
\(\sigma\mapsto(\sigma+\sigma^2)/2\), starting below one, decreases to
zero.  Given any desired positive exponent, choose the number of iterations
in advance, independently of \(L\).  There are then only finitely many
exceptional events to discard, each with probability tending to zero.
This proves \Cref{odd:all-power-geometry};
\Cref{odd:phase-clock} supplies the simultaneous phase estimates.
\end{proof}
